\section{Just enough category theory to be dangerous}
\begin{exer}[1.1.C.] Fix a field $k$. Let $(-)^{\vee\vee}$ be the double dual functor from the category of finite-dimensional $k$-vector spaces to itself. Show that $(-)^{\vee\vee}$ is naturally isomorphic to the identity functor.
    \begin{proof}
    For a fixed vector space $V$, define the map $\alpha_V: V \to V^{\vee\vee} $ by $v \mapsto (\phi\mapsto \phi(v))$. This is injective because if $\phi \mapsto \phi(v)$ is the zero map for every $\phi$, then $v=0$, so by finite-dimensionality $\alpha_V$ is an isomorphism. And it remains to check that for any map of vector spaces $f: V\to W$ and $v \in V$, that 
    \begin{align*}
        \alpha_W(f(v)) &= (\phi \mapsto \phi(f(v))) \\
        &= (\psi \mapsto \psi(v)) \circ f^\vee \\
        &= f^{\vee\vee}(\psi \mapsto \psi(v)) \\
        &= f^{\vee\vee}(\alpha_V(v))
    \end{align*}
    which completes the proof.
    \end{proof}
\end{exer}

\begin{exer}[1.1.D] Let $\mathcal{V}$ be the category whose objects are the $k$-vector spaces $k^n$ for $n\geq 0$, and whose morphisms are linear transformations. There is a functor $F$ from $\mathcal{V}$ to the category of f.d. $k$-vector spaces sending each vector space to itself. Show that $F$ is an equivalence of categories.
    \begin{proof}
    Following Vakil's remark about ignoring set-theoretic conditions, choose a basis for each finite dimensional $k$-vector space. For any vector space of the form $k^n$, choose the standard basis $\{e_i\}$.
    
    Then define a functor $G$ from f.d. $k$-vector spaces to $\mathcal{V}$ by sending an $n$-dimensional $k$-vector space to $k^n$. For a morphism $\phi$ from $V$ to an $m$-dimensional vector space $W$, let $\{v_1, ..., v_n\}$ be the chosen basis for $V$ and $\{w_1, ..., w_m\}$ the chosen basis for $W$. Then define $G\phi$ to be the matrix whose $(i,j)$th entry is the coefficient of $w_j$ in $\phi(v_i)$. 
    
    Define a natural transformation $\alpha$ from $FG$ to the identity functor where $\alpha_V: FGV \to V$ sends $FGV = k^n$ to $V$ by mapping the standard basis vector $e_i$ of $FGV$ to the chosen $i$th basis vector $v_i$ of $V$. Then $\alpha_V$ is an isomorphism. For naturality, we have for a map $\phi: V\to W$ that
    \begin{align*}
        \phi(\alpha_V(e_i)) &=\phi(v_i) \\
        &= \sum_j (G\phi)_{i,j} w_j \\
        &= \alpha_W(\sum_j(G\phi)_{i,j}e_j) \\
        &=\alpha_W (FG \phi)(e_i).
    \end{align*}
    In the other direction, $GF$ is the identity functor, which completes the proof.
    \end{proof}
\end{exer}

\begin{exer}[1.2.A] Show that any two initial objects are uniquely isomorphic, and any two final objects are uniquely isomorphic.
    \begin{proof}
    Let $i_1, i_2$ be initial. Then there is exactly one map $f_{12}: i_1 \to i_2$ and likewise $f_{21}: i_2 \to i_1$. Since there is also exactly one map from $i_1$ to itself, the identity, so because $f_{21}\circ f_{12}$ is one such map, it must be the identity. Likewise $f_{12}\circ f_{21}$ is the identity, which yields the unique isomorphisms. The argument is exactly the same for final objects.
    \end{proof}
\end{exer}
\begin{exer}[1.2.B]
What are the initial and final objects (if they exist) in $\catname{Ring}$, $\catname{Set}$, and $\catname{Top}$, or the category of subsets of a set under inclusion, or the category of open subsets of a topological space under inclusion?
\begin{soln}
In $\catname{Ring}$, the initial object is the integers, since $1$ mapping to $1$ dictates where the remaining integers go; the terminal object is the zero ring, since everything must map to zero. In $\catname{Set}$, the initial object is the empty set; the final object is the singleton. In $\catname{Top}$, the initial object is the empty space; the final object is the singleton. In the category of subsets or open sets, the initial object is the empty set/space; the terminal object is the entire set/space.
\end{soln}
\end{exer}

\begin{exer}[1.2.C]
    Let $A$ be a commutative ring and $S$ a multiplicative subset of $A$. Show that $f: A \mapsto S^{-1}A$ via $a \mapsto a/1$ is injective if and only if $S$ contains no zerodivisors.
    \begin{proof}
    First suppose $f$ is not injective, i.e. there is a nonzero element $a \in A$ such that $a/1 = 0/1$. So there exists some $s \in S$ such that $s(a1 - 01) = sa = 0$. Then $s$ is a zerodivisor. 

    Conversely, suppose $S$ contains a zerodivisor $s$, with $st = 0$ for $t$ nonzero. Then $0 = st = s(t1-01)$, $t/1 = 0/1 =0$, so $f(t) = t/1 = 0/1$ and $f$ fails to be injective.
    \end{proof}
\end{exer}

\begin{exer}[1.2.D] Verify that $\iota: A\to S^{-1}A$ is initial among $A$-algebras $B$ where every element of $S$ is sent to an invertible element in $B$, i.e. any map $\phi: A\to B$ sending all elements of $S$ to invertible elements in $B$ factors through $S^{-1}A$.
    \begin{proof}
        Define $\tilde{\phi}: S^{-1}A  \to B$ via $\tilde{\phi}(a/s) := \phi(a)\phi(s)^{-1}$. This makes sense since we have assumed that $\phi(s)$ is invertible. Let $a, b \in A$ and $s,t \in S$. Then 
        \[\tilde{\phi}(ab/st) = \phi(ab)\phi(st)^{-1} = \phi(a)\phi(b)\phi(s)^{-1}\phi(t)^{-1} = \tilde{\phi}(a/s) \tilde{\phi}(b/t)\]
        and 
        \begin{align*}
        \tilde{\phi}(a/s + b/t) &= \tilde{\phi}((at+sb)/st) \\
        &= \phi(at+sb)\phi(st)^{-1} \\
        &= \phi(a)\phi(t)/\phi(s)\phi(t)+\phi(s)\phi(b)/\phi(s)\phi(t) \\
        &= \tilde{\phi}(a/s) + \tilde{\phi}(b/t).
        \end{align*}
        so $\tilde{\phi}$ is a ring homomorphism. And $\tilde{\phi}(\iota(a)) = \tilde{\phi}(a/1) = \phi(a)/\phi(1) = \phi(a)$ as desired. 
    \end{proof}
\end{exer}

\begin{exer}[1.2.E] The localization of an $A$-module $M$ defined via the universal property, actually exists.
    \begin{proof}
        Following the hint, define $S^{-1}M$ to have elements of the form $m/s$ where $m \in M$ and $s \in S$, subject to the equivalence relation where $m/s \sim m'/t$ if there exists an element $s' \in S$ such that $s'(ms - m't) = 0$, with the standard operation and $S^{-1}A$-module operation. Define $\iota: M\to S^{-1}M$ via $m \mapsto m/1$. 

        Let $\phi: M \to N$ be an $A$-module map such that multiplication by any element of $s$ is an isomorphism $N\to N$. Then we can define a map $\tilde{\phi}: S^{-1}M\to N$ via $m/s \mapsto s^{-1}\phi(m)$, which makes sense taking $s^{-1}$ to be the inverse to multiplication by $s$. This is an $A$-module map and $\tilde{\phi}(\iota(m)) = \tilde{\phi}(a/1) = 1^{-1}\phi(a) = \phi(a)$.
    \end{proof}
\end{exer}

\begin{exer}[1.2.F] Properties of localization.
    \begin{enumerate}
        \item Localization commutes with finite products, i.e. for $A$-modules $M_1, ..., M_n$ provide an $A$-module isomorphism $S^{-1}(M_1\times \cdots M_n) \to S^{-1}M_1\times \cdots \times S^{-1}M_n$. 
        \begin{proof}
            Define our map via $(m_1, ..., m_n)/s\mapsto (m_1/s, ..., m_n/s)$. For surjectivity, given an element $(m_1/s_1, ..., m_n/s_n)$, the element
            \[(s_2\cdots s_nm_1, ..., s_1 \cdots s_{n-1}m_n)/s_1\cdots s_n\]
            will map to it. For injectivity, suppose $(m_1/s, ..., m_n/s)$ is zero, i.e. there exist elements $t_i \in S$ such that $t_im_i = 0$ for all $i$. Then  $t_1\cdots t_n(m_1, ..., m_n) = 0$, so $(m_1, ..., m_n)/s = 0$.
        \end{proof}
        \item Localization commutes with arbitrary direct sums.
        \begin{proof}
            Let $\Lambda$ be an indexing set. Define a map $S^{-1} (\bigoplus_{\lambda \in \Lambda} M_\lambda) \to \bigoplus_{\lambda \in \Lambda} S^{-1}M_{\lambda}$ via 
            \[(m_\lambda)_{\lambda \in \Lambda}/s \mapsto (m_\lambda/s)_{\lambda \in \Lambda}.\]
            For surjectivity, given an element $(m_\lambda / s_\lambda)_{\lambda \in \Lambda}$, we know that only finitely many indices are nonzero, say $\lambda_1, ..., \lambda_k$. Then let $s= \lambda_1\cdots \lambda_k$, and consider the element $(k_\lambda m_\lambda)_{\lambda \in \Lambda}$ where $k_{\lambda_i} = s/s_i$ for the nonzero indices and zero otherwise. Then it maps to what we want. For injectivity, suppose $(m_\lambda/s)_{\lambda \in \Lambda}$ is zero, i.e. there exist elements $t_\lambda$ such that $t_\lambda m_\lambda = 0$ for all $\lambda$. Again only finitely many of the $m_\lambda$ are nonzero, so take the corresponding $t_\lambda$ and let their product be $t$. Then $t(m_\lambda)_{\lambda \in \Lambda} = 0$ and we are done.
        \end{proof}
        \item Localization does not commute with infinite products, i.e. the same map defined as above is not always an isomorphism for infinite products. 
        \begin{proof}
            We know that $\Z$ localized at $\Z\setminus \{0\}$ is $\Q$, so that $\prod_{\N}\Q$ contains for example the element $x:= (1, 1/2, 1/3, ...)$. On the other hand, suppose there were an element in $(\Z\setminus \{0\})^{-1} (\Z^\N)$ mapping to $x$, i.e. some element $(a_1, a_2, ...)/s$ where $a_i \in \Z$ and $s \in \Z\setminus \{0\}$ such that $a_n/s = 1/n$ for all $n \in \N$. Then $na_n = s$ for all $n \in \N$, i.e. $s$ is an integer divisible by all natural numbers, a contradiction.
        \end{proof}
    \end{enumerate}
\end{exer}